\section{Práctica: Lab 2 — Sistema de detección de rostros}

El curso incluye una práctica (Lab 2) en la que se entrena un VAE con un conjunto de datos de rostros para descubrir automáticamente variables latentes (como el tono de piel, la pose, la iluminación, etc.) y utilizarlas para detectar sesgos y características subrepresentadas en el conjunto de datos.

\begin{figure}[H]
\centering
\includegraphics[width=0.85\textwidth]{slides-images/slide-89.jpg}
\caption{Lab 2: sistema de detección de rostros basado en VAE\protect\footnotemark}
\end{figure}
\footnotetext{Diapositivas, página 89.}

\begin{knowledgebox}{Sentido pedagógico del Lab 2}
Esta práctica combina teoría y práctica:
\begin{itemize}
    \item Construir un modelo VAE para procesar imágenes de rostros
    \item Descubrir características latentes en los datos mediante la perturbación de variables latentes
    \item Detectar sesgos en los datos de entrenamiento (por ejemplo, si ciertos tonos de piel o poses están subestimados)
    \item Utilizar modelos generativos para crear conjuntos de datos más justos y equilibrados
\end{itemize}
El código de la práctica está disponible en \href{https://github.com/MITDeepLearning/introtodeeplearning}{github.com/MITDeepLearning/introtodeeplearning}.
\end{knowledgebox}

\subsection{Resumen de la sección}

La práctica ayuda a los estudiantes a profundizar, mediante la experiencia directa, su comprensión de los VAE, en particular sobre cómo utilizar modelos generativos para descubrir y mitigar problemas de sesgo en los conjuntos de datos, un tema importante en la investigación sobre equidad en la IA.

